\SourceLeaf{310}%
in die Reihe
\[
\log y=\log x+\tfrac16xx+\tfrac{11}{180}x^4+\tfrac{191}{5670}x^6+\tfrac{2497}{113400}x^8+\text{u.\,s.\,w.}
\]
welche man leicht aus der bekannten
\[
y=x(1+\tfrac16xx+\tfrac3{40}x^4+\tfrac5{112}x^6+\tfrac{35}{1152}x^8+\text{u.\,s.\,w.})
\]
ableitet. Wendet man dieselbe zuvörderst an auf die Gleichung
\[
\operatorname{tang}B\mathord{.}\sin T\mathord{.}\operatorname{tang}\tfrac12r=\sin\tfrac12t
\]
oder
\[
\tfrac12m\tau=\sin\tfrac12t
\]
indem man \(x=\frac12m\tau\), \(y=\frac12t\) setzt, so wird (I)
\[
\begin{aligned}
\log t={}&\log\tau+\log m+\tfrac1{24}mm\tau\tau+\tfrac{11}{2880}m^4\tau^4\\
&+\tfrac{191}{362880}m^6\tau^6+\tfrac{2497}{29030400}m^8\tau^8+\text{u.\,s.\,w.}
\end{aligned}
\]

\indent Eben so, aus der Anwendung auf die Gleichung
\[
\frac{\sin T\sin\frac12r}{\cos B}=\sin\tfrac12l
\]
oder
\[
\tfrac12n\lambda=\sin\tfrac12l
\]
ergibt sich (II)
\[
\begin{aligned}
\log l={}&\log\lambda+\log n+\tfrac1{24}nn\lambda\lambda+\tfrac{11}{2880}n^4\lambda^4\\
&+\tfrac{191}{362880}n^6\lambda^6+\tfrac{2497}{29030400}n^8\lambda^8+\text{u.\,s.\,w.}
\end{aligned}
\]

\indent Die dritte Anwendung wird gemacht auf die Gleichung
\[
\frac{\cos B\operatorname{tang}\frac12l}{\operatorname{tang}T}=\sin\tfrac12b
\]
nachdem derselben vermöge der Substitutionen
\[
\begin{aligned}
\operatorname{tang}\tfrac12l&=\frac{n\lambda}{2\BareRoot{(1-\frac14nn\lambda\lambda)}}\\
\frac{\cos B}{\operatorname{tang}T}&=\frac\beta\lambda
\end{aligned}
\]
folgende Gestalt gegeben ist
\[
\frac{n\beta}{2\BareRoot{(1-\frac14nn\lambda\lambda)}}=\sin\tfrac12b
\]
Es ergibt sich dann (III)
